\documentclass[11pt]{article}
\usepackage[a4paper,margin=22mm]{geometry}
\usepackage{fontspec}
\setmainfont{DejaVu Serif}
\setsansfont{DejaVu Sans}
\usepackage{microtype}
\usepackage{xcolor}
\usepackage{hyperref}
\usepackage{enumitem}
\usepackage{parskip}
\definecolor{Accent}{HTML}{971D32}
\definecolor{Muted}{HTML}{666666}
\hypersetup{colorlinks=true,urlcolor=Accent,linkcolor=Accent}
\setlist{leftmargin=1.6em,itemsep=0.3em,topsep=0.35em}
\setcounter{tocdepth}{2}
\renewcommand{\contentsname}{Contents}
\newcommand{\sectionrule}{\par\vspace{-0.5em}\noindent\textcolor{Accent}{\rule{\linewidth}{0.6pt}}\par}
\begin{document}

\begin{center}
{\sffamily\bfseries\color{Accent} TOEFL WORD OF THE DAY\par}
\vspace{8mm}
{\fontsize{42}{48}\selectfont\bfseries dovetail\par}
\vspace{2mm}
{\Large\sffamily\color{Accent} /ˈdʌvˌteɪl/\par}
\vspace{2mm}
{\small\sffamily Local audio phrase: dovetail\par}
{\small\sffamily Audio file: audios/dovetail.mp3\par}
\vspace{2mm}
{\large\sffamily\color{Muted} intransitive and transitive verb; countable noun\par}
\vspace{5mm}
{\sffamily fit together closely \textbullet\ a tightly fitting woodworking joint\par}
\vspace{6mm}
{\large Dovetail usually means that two ideas, plans, or pieces of evidence fit together very well; it also names the woodworking joint behind this figurative meaning.\par}
\vspace{6mm}
\begin{minipage}{0.84\textwidth}
\itshape “Her explanation dovetails neatly with the evidence presented earlier.”
\end{minipage}\par
\vspace{10mm}
\textbf{Date:} August 19th 2026\quad\textbullet\quad \textbf{Level:} B1--mid-B2 TOEFL
\vfill
Created and compiled by \textbf{Amir Kooshky}\par
\vspace{2mm}
\href{https://t.me/KooshkyTOEFL}{Telegram Channel}\quad\textbullet\quad
\href{https://www.instagram.com/kooshkytoefl}{Instagram}\quad\textbullet\quad
\href{https://t.me/KooshkyTOEFL\_pv}{Personal Telegram}
\end{center}
\newpage
\tableofcontents
\newpage

\section{Meanings and usage}
\sectionrule
\subsection{Meaning 1: Fit together closely}
\textbf{Part of speech:} intransitive and transitive verb

\textbf{IPA:} /ˈdʌvˌteɪl/

\textbf{Pronunciation phrase:} dovetail

\textbf{Local audio file:} audios/dovetail.mp3

\textbf{Register:} neutral to formal; common in academic, professional, organizational, and analytical writing

\textbf{Definition:} to fit together closely and successfully, or to make two things fit together in this way

\textbf{In simpler words:} fit together very well

\textbf{Typical pattern:} dovetail with + idea, evidence, plan, schedule, or goal

\subsubsection{Natural examples}
\begin{enumerate}
  \item The results of the two studies dovetail with each other and support the same conclusion.
  \item The new training program was designed to dovetail with employees' existing responsibilities.
  \item Her explanation dovetails neatly with the evidence presented earlier.
  \item The historical records dovetail with archaeological evidence from the same period.
  \item The two departments adjusted their schedules so that their activities would dovetail.
  \item The author's economic argument does not fully dovetail with the social evidence discussed later.
  \item The researchers tried to dovetail their different methods into a single investigation.
  \item The proposed reforms dovetail closely with the university's long-term goals.
\end{enumerate}

\subsubsection{Nuance and implication}
\textbf{Central nuance:} Dovetail suggests more than simple similarity. The different parts complement one another and fit together in a useful or logical way.

\textbf{Usual implication:} The things being compared often support, reinforce, or coordinate with one another.

\textbf{Compared with conflict:} Things that dovetail fit together successfully, while things that conflict are inconsistent or interfere with one another.

\subsubsection{Grammar notes}
\textbf{How to use it:} It is often used without a direct object and followed by with, as in The findings dovetail with earlier research. It can also be used with an object when someone deliberately combines things.

\textbf{What can follow it:} Common subjects and objects include findings, evidence, theories, plans, schedules, goals, policies, and activities.

\textbf{Common preposition:} Use dovetail with when saying that one thing fits well with another.

\textbf{Position in a sentence:} Adverbs such as neatly, closely, perfectly, and well often appear with dovetail.

\subsubsection{Useful patterns}
\textbf{Pattern:} dovetail with + noun

\textbf{Use:} Use this when two ideas, findings, plans, or activities fit together well.

\textbf{Example:} The new evidence dovetails with the earlier findings.

\textbf{Pattern:} dovetail neatly or closely with + noun

\textbf{Use:} Use an adverb to emphasize how well two things fit together.

\textbf{Example:} The policy dovetails closely with the government's environmental goals.

\textbf{Pattern:} dovetail + activities, plans, or methods

\textbf{Use:} Use this when people deliberately coordinate different things.

\textbf{Example:} The two teams dovetailed their research activities to avoid unnecessary duplication.

\textbf{Pattern:} dovetail into + larger whole

\textbf{Use:} Use this when separate parts combine smoothly into a broader plan or system.

\textbf{Example:} The individual projects dovetail into a broader regional development strategy.

\textbf{Pattern:} fail to dovetail with + noun

\textbf{Use:} Use this when two ideas or plans do not fit together as expected.

\textbf{Example:} The theory fails to dovetail with several important observations.

\subsubsection{Common wrong usage}
\paragraph{Correction 1}
\textbf{Wrong:} The new findings dovetail to the previous research.

\textbf{Correct:} The new findings dovetail with the previous research.

\textbf{Why:} Use dovetail with when two things fit together well.

\paragraph{Correction 2}
\textbf{Wrong:} The two theories are dovetail each other.

\textbf{Correct:} The two theories dovetail with each other.

\textbf{Why:} Use dovetail with each other rather than be dovetail each other.

\paragraph{Correction 3}
\textbf{Wrong:} The evidence dovetails similarly with the theory because both are identical.

\textbf{Correct:} The evidence dovetails with the theory because the two fit together well.

\textbf{Why:} Dovetail does not mean identical. It means that different things fit together successfully.

\paragraph{Correction 4}
\textbf{Wrong:} The manager dovetailed with the two schedules.

\textbf{Correct:} The manager dovetailed the two schedules.

\textbf{Why:} When a person deliberately coordinates two things, those things can be direct objects of dovetail.

\subsection{Meaning 2: Tightly fitting woodworking joint}
\textbf{Part of speech:} countable noun

\textbf{IPA:} /ˈdʌvˌteɪl/

\textbf{Pronunciation phrase:} dovetail

\textbf{Local audio file:} audios/dovetail.mp3

\textbf{Register:} neutral; mainly used in woodworking, craftsmanship, archaeology, and descriptions of furniture construction

\textbf{Definition:} a type of joint in which projecting pieces of wood fit tightly into matching spaces in another piece

\textbf{In simpler words:} a strong interlocking wood joint

\textbf{Typical pattern:} cut, make, or use + dovetails

\subsubsection{Natural examples}
\begin{enumerate}
  \item The carpenter used dovetails to join the sides of the wooden drawer.
  \item Well-made dovetails can create a strong joint without relying entirely on nails or screws.
  \item Archaeologists found furniture fragments containing carefully cut dovetails.
  \item Traditional cabinetmakers often used dovetails because the joint resists being pulled apart.
  \item The wooden box was assembled with several hand-cut dovetails.
  \item The shape of each dovetail allows two pieces of wood to lock together firmly.
  \item Students in the woodworking class practiced cutting dovetails by hand.
\end{enumerate}

\subsubsection{Nuance and implication}
\textbf{Central nuance:} The projecting part is wider at the end than at its base, creating an interlocking shape that resists being pulled apart.

\textbf{Usual implication:} The physical way the pieces fit together gave rise to the figurative verb meaning of fitting together successfully.

\subsubsection{Grammar notes}
\textbf{How to use it:} Use a dovetail for one joint or projecting piece and dovetails for more than one.

\textbf{What can follow it:} It often appears with words such as joint, drawer, cabinet, wood, cut, and hand-cut.

\textbf{Noun use:} This is a countable noun, so the singular normally needs a, an, or the.

\subsubsection{Useful patterns}
\textbf{Pattern:} a dovetail joint

\textbf{Use:} Use this for the complete interlocking woodworking joint.

\textbf{Example:} The cabinet is held together with traditional dovetail joints.

\textbf{Pattern:} cut dovetails

\textbf{Use:} Use this for making the interlocking shapes in wood.

\textbf{Example:} Experienced woodworkers can cut dovetails by hand.

\textbf{Pattern:} hand-cut dovetails

\textbf{Use:} Use this for dovetails made manually rather than by a machine.

\textbf{Example:} The antique drawer contains finely made hand-cut dovetails.

\subsubsection{Common wrong usage}
\paragraph{Correction 1}
\textbf{Wrong:} The carpenter made a dovetail with two boards together.

\textbf{Correct:} The carpenter joined the two boards with a dovetail joint.

\textbf{Why:} Use with a dovetail joint when naming the type of joint used to connect pieces of wood.

\paragraph{Correction 2}
\textbf{Wrong:} The drawer has dovetail to make it stronger.

\textbf{Correct:} The drawer has dovetails to make it stronger.

\textbf{Why:} Dovetail is countable, so use the plural when referring to several joints or interlocking parts.

\section{Word family}
\sectionrule
\subsection{dovetailed}
\textbf{Part of speech:} adjective

\textbf{IPA:} /ˈdʌvˌteɪld/

\textbf{Definition:} joined or coordinated so that different parts fit together closely

\textbf{Relationship to the headword:} This adjective describes things that have been physically joined or carefully coordinated so that they fit together.

\textbf{Grammar pattern:} dovetailed + activities, plans, joints, or corners

\textbf{Usage note:} It is less common than the verb dovetail.

\subsubsection{Examples}
\begin{enumerate}
  \item The project depended on several carefully dovetailed activities.
  \item The cabinet has dovetailed corners that provide additional strength.
\end{enumerate}

\section{Common collocations}
\sectionrule
\subsection{dovetail with}
\textbf{Applies to:} Fit together closely

\textbf{Meaning:} fit or agree closely with something

\textbf{Pattern:} dovetail with + evidence, finding, plan, goal, or theory

\textbf{Example:} The findings dovetail with previous research on the same population.

\subsection{dovetail neatly with}
\textbf{Applies to:} Fit together closely

\textbf{Meaning:} fit together in a particularly smooth or logical way

\textbf{Example:} The new explanation dovetails neatly with the archaeological evidence.

\subsection{dovetail closely with}
\textbf{Applies to:} Fit together closely

\textbf{Meaning:} correspond or coordinate very well with something

\textbf{Example:} The course objectives dovetail closely with the department's broader goals.

\subsection{dovetail perfectly with}
\textbf{Applies to:} Fit together closely

\textbf{Meaning:} fit together extremely well

\textbf{Example:} The revised schedule dovetails perfectly with the transportation timetable.

\subsection{dovetail with existing plans}
\textbf{Applies to:} Fit together closely

\textbf{Meaning:} coordinate successfully with plans that already exist

\textbf{Example:} The new initiative was designed to dovetail with existing development plans.

\subsection{dovetail with the evidence}
\textbf{Applies to:} Fit together closely

\textbf{Meaning:} be consistent with and supported by available evidence

\textbf{Example:} Her account dovetails with the physical evidence collected at the site.

\subsection{dovetail joint}
\textbf{Applies to:} Tightly fitting woodworking joint

\textbf{Meaning:} a woodworking joint formed by interlocking wedge-shaped projections

\textbf{Example:} A dovetail joint connects the side of the drawer to its front.

\subsection{hand-cut dovetails}
\textbf{Applies to:} Tightly fitting woodworking joint

\textbf{Meaning:} dovetail joints shaped manually rather than by machine

\textbf{Example:} The antique cabinet contains carefully made hand-cut dovetails.

\section{Synonyms and antonyms}
\sectionrule
\subsection{Close synonyms}
\subsubsection{correspond}
\textbf{Part of speech:} intransitive verb

\textbf{Applies to:} Fit together closely

\textbf{Shared meaning:} Both can describe agreement between findings, ideas, or information.

\textbf{Difference:} Correspond means match or be similar in a relevant way. Dovetail more strongly suggests that separate things complement each other and fit together into a coherent whole.

\textbf{Example:} The survey results correspond closely with earlier estimates.

\subsubsection{align}
\textbf{Part of speech:} intransitive or transitive verb

\textbf{Applies to:} Fit together closely

\textbf{Shared meaning:} Both can describe plans or goals that work well together.

\textbf{Difference:} Align often means bring ideas, goals, or activities into agreement or the same direction. Dovetail suggests that different parts complement each other and fit together closely.

\textbf{Example:} The new policy aligns with the organization's long-term goals.

\subsubsection{mesh}
\textbf{Part of speech:} intransitive verb

\textbf{Applies to:} Fit together closely

\textbf{Shared meaning:} Both can mean that different parts work successfully together.

\textbf{Difference:} Mesh is an informal to neutral synonym meaning work or fit together successfully. Dovetail is often more formal and suggests a particularly neat or complementary fit.

\textbf{Example:} The two teaching methods mesh well together.

\subsubsection{integrate}
\textbf{Part of speech:} transitive or intransitive verb

\textbf{Applies to:} Fit together closely

\textbf{Shared meaning:} Both can describe combining separate elements successfully.

\textbf{Difference:} Integrate means combine parts into a larger whole. Dovetail emphasizes how well those different parts fit or coordinate with one another.

\textbf{Example:} The researchers integrated several sources of data into one model.

\subsection{Direct or useful antonyms}
\subsubsection{conflict}
\textbf{Part of speech:} intransitive verb

\textbf{Applies to:} Fit together closely

\textbf{Opposition:} Things that conflict are inconsistent or interfere with one another, while things that dovetail fit together successfully.

\textbf{Example:} The witness's account conflicts with the physical evidence.

\subsubsection{clash}
\textbf{Part of speech:} intransitive verb

\textbf{Applies to:} Fit together closely

\textbf{Opposition:} Things that clash do not combine or work well together, while things that dovetail complement one another.

\textbf{Example:} The two departments' schedules clash on several important dates.

\section{Words learners may confuse}
\sectionrule
\subsection{overlap}
\textbf{Part of speech:} verb

\textbf{Why learners confuse them:} Both can describe relationships between activities, ideas, or areas of study, but overlap emphasizes shared content rather than complementary fit.

\textbf{Key difference:} Overlap means share some of the same area, time, features, or content. Dovetail means fit together successfully, often because different parts complement one another.

\textbf{dovetail:} The two research methods dovetail well because each provides information the other lacks.

\textbf{overlap:} The two courses overlap because they cover several of the same topics.

\subsection{align}
\textbf{Part of speech:} verb

\textbf{Why learners confuse them:} Both are frequently used when discussing ideas, plans, goals, and organizational activities that work well together.

\textbf{Key difference:} Align emphasizes agreement in direction, position, or goals. Dovetail emphasizes different parts fitting together neatly and complementing one another.

\textbf{dovetail:} The new evidence dovetails with earlier findings.

\textbf{align:} The department changed its priorities to align with the university's strategy.

\section{Etymology}
\sectionrule
\textbf{Origin:} Dovetail comes from the woodworking joint whose projecting pieces resemble the spread tail of a dove.

\textbf{Development:} Because the pieces of a dovetail joint fit tightly together, the word later developed the figurative meaning of ideas, plans, or activities fitting together successfully.

\textbf{Memory link:} Imagine two pieces of wood locking neatly together: ideas that dovetail also fit together neatly.

\section{Practice exercises}
\sectionrule
\subsection{Word family choice}
Choose the best member of the word family for each blank.

\begin{enumerate}
  \item The project involved several carefully \_\_\_\_\_ activities that had to be completed in sequence.
  \item The new findings \_\_\_\_\_ closely with the results of the earlier study.
\end{enumerate}

\subsection{Correct or incorrect?}
Decide whether each sentence is correct. Correct the inaccurate ones.

\begin{enumerate}
  \item The new evidence dovetails with the earlier findings.
  \item The new policy dovetails to the university's existing plans.
  \item The two theories dovetail because they are completely identical.
  \item The researchers dovetailed their methods into a single study.
\end{enumerate}

\subsection{Collocation practice}
Complete each sentence with a natural collocation.

\begin{enumerate}
  \item The new findings dovetail \_\_\_\_\_ earlier research on the same topic.
  \item The training program was designed to dovetail \_\_\_\_\_ employees' existing responsibilities.
  \item Her explanation dovetails \_\_\_\_\_ with the available evidence.
  \item The antique drawer was constructed using hand-cut \_\_\_\_\_.
\end{enumerate}

\subsection{Complete answer key}
\subsubsection{Word family choice}
\begin{enumerate}
  \item \textbf{dovetailed.} The project involved several carefully dovetailed activities that had to be completed in sequence. \emph{Use the adjective dovetailed to describe activities coordinated so that they fit together.}
  \item \textbf{dovetail.} The new findings dovetail closely with the results of the earlier study. \emph{Use the verb dovetail when two sets of findings fit together well.}
\end{enumerate}

\subsubsection{Correct or incorrect?}
\begin{enumerate}
  \item \textbf{Correct} \emph{Dovetail with is the standard pattern when two pieces of evidence fit together well.}
  \item \textbf{Incorrect}. The new policy dovetails with the university's existing plans. \emph{Use dovetail with, not dovetail to.}
  \item \textbf{Incorrect}. The two theories dovetail because their different explanations complement each other. \emph{Dovetail means fit together well, not necessarily be identical.}
  \item \textbf{Correct} \emph{Dovetail can be transitive when someone deliberately coordinates or combines different things.}
\end{enumerate}

\subsubsection{Collocation practice}
\begin{enumerate}
  \item \textbf{with.} The new findings dovetail with earlier research on the same topic. \emph{Use dovetail with when two things fit together successfully.}
  \item \textbf{with.} The training program was designed to dovetail with employees' existing responsibilities. \emph{Dovetail with means coordinate well with something else.}
  \item \textbf{neatly.} Her explanation dovetails neatly with the available evidence. \emph{Dovetail neatly with emphasizes a particularly good fit.}
  \item \textbf{dovetails.} The antique drawer was constructed using hand-cut dovetails. \emph{Dovetails are the interlocking parts used in a traditional woodworking joint.}
\end{enumerate}

\vfill
\begin{center}
\small\color{Muted} Created and compiled by \textbf{Amir Kooshky}\\
\href{https://t.me/KooshkyTOEFL}{Telegram Channel}\quad\textbullet\quad
\href{https://www.instagram.com/kooshkytoefl}{Instagram}\quad\textbullet\quad
\href{https://t.me/KooshkyTOEFL\_pv}{Personal Telegram}
\end{center}
\end{document}
